\documentclass{relatorio}
\usepackage[utf8]{inputenc}
\usepackage[T1]{fontenc}
\usepackage{geometry}
\usepackage{graphicx}
\usepackage{float}
\usepackage{hyperref}
\usepackage[english,indonesian]{babel}

\title{Optimasi Pemrosesan Template LaTeX Menggunakan Pendekatan Berbasis AI}
\author{Academic Support Agent}
\date{\today}

\begin{document}
\maketitle

\begin{abstract}

\end{abstract}

\chapter{PENDAHULUAN}


\chapter{TINJAUAN PUSTAKA}


% Original content (commented out):
% % !TeX program = xelatex
% !TeX encoding = UTF-8 Unicode
% !BIB program = biber

% essa classe aceita as seguintes opções:
% abnd ou ieee pra estilo de citaçõ
% english ou french pra língua dos slides (não passar nada é português)
% minted pra usar o pacote minted (configurado para python e C)
\documentclass[minted]{relatorio}

% Todo comando \add trabalha com uma lista e todo comando \set trabalha com um
% único elemento, ou seja, pode haver apenas um \setorientador mas vários
% \addauthor. Todo comando \add e \set tem um \get correspondente (por exemplo,
% \getauthors, \getsubject). Esses comandos são definidos na classe
% relatorio.cls. No caso dos \set, os \get são implicitos, e não você não vai
% encontrá-los na classe.

\title{Exemplo de Relatório}
\addauthor{Álan Crístoffer e Sousa}
\addauthor{Lanna Gontijo Ferreira}
\setsubject{Controle Digital e Laboratório}

\preamble{}

% Essa classe usa o biblatex com backend biber ao invés do bibtex.
\addbibresource{bibliothek.bib}

% Controla o texto das captions das figuras. Neste caso, centraliza.
\captionsetup{justification=centering}

\begin{document}
\input{abstract}

\maketitle{}

\input{section1}
\input{section2}
\input{section3}
\input{section4}
\input{section5}

\nocite{*}
\printbibliography{}
\end{document}


\end{document}